% abstract.tex
\begin{abstract}
Detecting unseen multi-turn jailbreaks requires domain adaptation, but the form of adaptation matters.
Multi-turn jailbreak attacks spread adversarial intent across benign-seeming turns, and existing defenses fail to generalize across attack families: Llama Guard~3, a widely deployed safety classifier, detects only 1.8\% of multi-turn jailbreaks in our evaluation.
We isolate the factors behind cross-attack generalization on synthetic multi-turn data through controlled ablation: thirteen encoder configurations at three scales (125 to 355M parameters) decompose the contributions of domain exposure, classification supervision, label quality, and granularity on three out-of-distribution (OOD) attack families and two training compositions.
On single-source synthetic data, continued pretraining closes roughly 60 percentage points of the OOD gap without labels. This effectiveness is condition-dependent: on mixed-LLM training data, unsupervised pretraining collapses while classification-based adaptation remains the strongest variant (DD OOD F1\,=\,0.818 vs.\ 0.291 without adaptation).
Across three architectures (125--355M), the set of attacks requiring adaptation shrinks with scale, though architecture and pretraining differences prevent isolating scale as the sole factor.
Classification supervision is associated with substantially lower prediction variance on synthetic evaluations (observed across 3 seeds), with limited sensitivity to label correctness, but this benefit reverses on human-authored data.
These findings suggest classification DAPT provides robustness through optimization regularization, not semantic feature learning\footnote{Code and data are available at \url{https://anonymous.4open.science/r/persuasion_turn_jailbreak_detect-C207/}.}.
\end{abstract}
